% J2.7: ajustar, diagnosticar, mejorar y usar un modelo de regresión.
\documentclass[tikz,border=4pt]{standalone}
\usepackage{resonancia}
\begin{document}
\begin{tikzpicture}
  \node[cajaAzul, minimum width=2.4cm] (a) at (0,0) {\textbf{1. Ajustar}\\[-1pt]{\footnotesize \texttt{lm(@formula(f\,\~{}\,T))}}};
  \node[cajaVerde, minimum width=2.4cm] (d) at (3.5,0) {\textbf{2. Diagnosticar}\\[-1pt]{\footnotesize residuos, QQ, influencia}};
  \node[cajaNaranja, minimum width=2.4cm] (m) at (7.0,0) {\textbf{3. Mejorar}\\[-1pt]{\footnotesize términos, AIC, \texttt{ftest}}};
  \node[caja, draw=rVioleta, fill=rVioleta!8, minimum width=2.4cm] (u) at (10.5,0) {\textbf{4. Usar}\\[-1pt]{\footnotesize predecir, compensar}};
  \draw[flecha] (a) -- (d); \draw[flecha] (d) -- (m); \draw[flecha] (m) -- (u);
  \draw[flecha, color=rGris, dashed] (m.south) -- ++(0,-0.7) -| (d.south);
  \node[etiqueta, color=rGris, anchor=north] at (5.2,-0.95) {se itera entre 2 y 3 hasta que los residuos no tengan estructura};
\end{tikzpicture}
\end{document}
